\clearpage\thispagestyle{empty}
\pdfbookmark[0]{Quellen zur Besprechung der Connaissance des tems für 1829}{connaissance1829-sources}
\begin{center}\Large\bfseries Quellen zur Besprechung der Connaissance des tems für 1829\end{center}

Die Übertragung folgt Carl Friedrich Gauß’ Besprechung in den \emph{Werken}, Bd.~VI (Göttingen, 1874), S.~645--648. Der Erstabdruck erschien in den \emph{Göttingischen gelehrten Anzeigen}, Stück~6, 10.~Januar 1828, S.~49--56. Das Bureau des longitudes gab die \emph{Connaissance des tems} für 1829 im Jahr 1826 bei Bachelier in Paris heraus.

KI-Übertragung aus dem historischen Deutsch in modernes Deutsch. KI-erstellter Satz.

\href{https://gdz.sub.uni-goettingen.de/id/PPN236007467}{Werkausgabe}. Erstdruck, S.~\href{https://images.p.sub.uni-goettingen.de/iiif/image/gdz:PPN385030444_1828_1:00000051/full/1650,/0/default.jpg}{49}, \href{https://images.p.sub.uni-goettingen.de/iiif/image/gdz:PPN385030444_1828_1:00000052/full/1650,/0/default.jpg}{50}, \href{https://images.p.sub.uni-goettingen.de/iiif/image/gdz:PPN385030444_1828_1:00000053/full/1650,/0/default.jpg}{51}, \href{https://images.p.sub.uni-goettingen.de/iiif/image/gdz:PPN385030444_1828_1:00000054/full/1650,/0/default.jpg}{52}, \href{https://images.p.sub.uni-goettingen.de/iiif/image/gdz:PPN385030444_1828_1:00000055/full/1650,/0/default.jpg}{53}, \href{https://images.p.sub.uni-goettingen.de/iiif/image/gdz:PPN385030444_1828_1:00000056/full/1650,/0/default.jpg}{54}, \href{https://images.p.sub.uni-goettingen.de/iiif/image/gdz:PPN385030444_1828_1:00000057/full/1650,/0/default.jpg}{55}, \href{https://images.p.sub.uni-goettingen.de/iiif/image/gdz:PPN385030444_1828_1:00000058/full/1650,/0/default.jpg}{56}. \href{https://cdt.imcce.fr/items/show/178}{IMCCE-Katalog zum Jahrgang 1829}.

\clearpage
